% ==========================================
% Week 8: SYSTEM, CONTROL VOLUME, RTT
% ==========================================
\section{Week 8 - System, Control Volume, and Reynolds Transport Theorem}

This section closes the chapter sequence by linking the Lagrangian and Eulerian viewpoints and showing how conservation laws are written for control volumes.

\subsection{System vs. Control Volume}

\begin{itemize}
    \item \textbf{System:} A fixed amount of mass. Its boundary may move and deform with the material. This is the natural \textbf{Lagrangian} viewpoint.
    \item \textbf{Control Volume (CV):} A region in space through which mass may enter or leave. This is the natural \textbf{Eulerian} viewpoint.
\end{itemize}

The purpose of the Reynolds Transport Theorem is to connect these two descriptions.

\subsection{Reynolds Transport Theorem (RTT)}

\textbf{Formal derivation:} Munson 4.4.1

Consider a system at time $t=t_0$ such that the control volume and the system coincide:
\begin{equation}
    CV \text{ at } t=t_0 = SYS
\end{equation}

Let
\begin{itemize}
    \item $B$ be an extensive quantity
    \item $b$ be the corresponding intensive quantity defined by
    \begin{equation}
        b=\frac{B}{m}
    \end{equation}
\end{itemize}

The amount of $B$ inside the control volume and inside the system can be written as
\begin{equation}
    B_{CV}=\int_{CV}\rho b\,dV,
    \qquad
    B_{SYS}=\int_{SYS}\rho b\,dV
\end{equation}

The rate of change of $B$ for the system is related to the control-volume description by
\begin{equation}
    \frac{D B_{SYS}}{Dt}
    =
    \frac{\partial B_{CV}}{\partial t}
    +
    \dot{B}_{out}
    -
    \dot{B}_{in}
\end{equation}

The flux terms are
\begin{itemize}
    \item Flux of $B$ into the control volume:
    \begin{equation}
        \dot{B}_{in}=b_{in}\dot{m}_{in}=b_{in}\rho_{in}A_{in}V_{in}
    \end{equation}

    \item Flux of $B$ out of the control volume:
    \begin{equation}
        \dot{B}_{out}=b_{out}\dot{m}_{out}=b_{out}\rho_{out}A_{out}V_{out}
    \end{equation}
\end{itemize}

\begin{mainbox}{Simple Version of RTT}
\begin{equation*}
    \left(\frac{DB}{Dt}\right)_{SYS}
    =
    \left(\frac{\partial B}{\partial t}\right)_{CV}
    +
    (b\rho VA)_{out}
    -
    (b\rho VA)_{in}
\end{equation*}
\end{mainbox}

\subsubsection{General Form of RTT}

For a general surface, the mass flow rate through an area element is
\begin{equation}
    \dot{m}=\int_A \rho \vec{V}\cdot\hat{n}\,dA
\end{equation}

Therefore, the net rate at which the property $B$ crosses the control surface is
\begin{equation}
    \dot{B}
    =
    \int_{CS} b\rho V_n\,dA
    =
    \dot{B}_{out}-\dot{B}_{in}
\end{equation}
where
\begin{itemize}
    \item $\hat{n}$ is the outward unit normal vector
    \item $dA$ is an area element on the control surface
    \item $dV$ is a volume element in the control volume
    \item $V_n=\vec{V}\cdot\hat{n}$ is the normal component of velocity
\end{itemize}

Since $CV=SYS$ at the instant $t=t_0$, the general form of the Reynolds Transport Theorem is
\begin{mainbox}{Reynolds Transport Theorem (General Form)}
\begin{equation*}
    \underbrace{\left(\frac{DB}{Dt}\right)_{SYS}}_{\text{Lagrangian system}}
    =
    \underbrace{\frac{\partial}{\partial t}\int_{CV} b\rho\,dV}_{\text{Eulerian volume contribution}}
    +
    \underbrace{\int_{CS} b\rho(\vec{V}\cdot\hat{n})\,dA}_{\text{Eulerian surface contribution}}
\end{equation*}
\end{mainbox}

\textbf{Notes:}
\begin{itemize}
    \item For \textbf{steady flow},
    \begin{equation}
        \frac{\partial}{\partial t}\int_{CV} b\rho\,dV = 0
    \end{equation}
    so that
    \begin{equation}
        \frac{DB_{SYS}}{Dt}
        =
        \int_{CS} b\rho(\vec{V}\cdot\hat{n})\,dA
    \end{equation}

    \item For a \textbf{moving, non-deforming control volume}, one uses the relative velocity
    \begin{equation}
        \vec{W}=\vec{V}-\vec{V}_{CV}
    \end{equation}

    \item In practical calculations, the closed-surface integral is often split into separate inlet and outlet contributions.
\end{itemize}

\subsection{Control Volume Approach to Conservation Laws}

Using RTT, the main Lagrangian conservation laws can be rewritten in Eulerian control-volume form:
\begin{enumerate}
    \item Mass conservation
    \item Newton's second law
    \item Energy conservation (First Law of Thermodynamics)
\end{enumerate}

\subsection{Mass Conservation}

For mass conservation, the extensive quantity is the mass itself:
\begin{equation}
    B=M
\end{equation}
Hence the associated intensive quantity is
\begin{equation}
    b=1
\end{equation}

Substituting into the general RTT gives
\begin{equation}
    \left(\frac{DM}{Dt}\right)_{SYS}
    =
    \frac{\partial}{\partial t}\int_{CV}\rho\,dV
    +
    \int_{CS}\rho(\vec{V}\cdot\hat{n})\,dA
\end{equation}

In words:
\begin{quote}
Rate of change of the mass of the coincident system
=
rate of change of mass inside the control volume
+
net mass flow rate through the control surface.
\end{quote}

For a closed system, mass is conserved, so
\begin{equation}
    \left(\frac{DM}{Dt}\right)_{SYS}=0
\end{equation}

This gives the integral continuity equation:
\begin{eqbox}{Continuity Equation}
\begin{equation*}
    0
    =
    \frac{\partial}{\partial t}\int_{CV}\rho\,dV
    +
    \int_{CS}\rho(\vec{V}\cdot\hat{n})\,dA
\end{equation*}
\end{eqbox}

\subsubsection{Example: Nozzle (Stationary Flow)}

Consider a nozzle with:
\begin{itemize}
    \item \textbf{Inlet (Area $A_0$):} uniform velocity profile $V_0$, radius $R_0$
    \item \textbf{Outlet (Area $A_1$):} parabolic velocity profile
    \begin{equation}
        V(r)=V_m\left(1-\frac{r^2}{R_1^2}\right)
    \end{equation}
    with outlet radius $R_1$
\end{itemize}

\textbf{Assumptions:}
\begin{itemize}
    \item Steady flow
    \item Constant density, $\rho=\text{const}$
\end{itemize}

Apply the continuity equation:
\begin{equation}
    \frac{\partial}{\partial t}\int_{CV}\rho\,dV
    +
    \int_{CS}\rho(\vec{V}\cdot\hat{n})\,dA
    =
    0
\end{equation}

Because the flow is steady, the accumulation term vanishes:
\begin{equation}
    \frac{\partial}{\partial t}\int_{CV}\rho\,dV=0
\end{equation}
Therefore,
\begin{equation}
    \int_{CS}\rho(\vec{V}\cdot\hat{n})\,dA=0
\end{equation}
and since $\rho$ is constant,
\begin{equation}
    \int_{CS}(\vec{V}\cdot\hat{n})\,dA=0
\end{equation}

Split the control-surface integral into inlet and outlet contributions:
\begin{equation}
    \int_{A_0}(\vec{V}\cdot\hat{n})\,dA
    +
    \int_{A_1}(\vec{V}\cdot\hat{n})\,dA
    =
    0
\end{equation}

\textbf{Inlet contribution:} \\
At the inlet, the outward normal points opposite to the incoming flow, so
\begin{equation}
    \vec{V}\cdot\hat{n}=-V_0
\end{equation}
Hence,
\begin{equation}
    \int_{A_0}(\vec{V}\cdot\hat{n})\,dA
    =
    -V_0\int_{A_0}dA
    =
    -V_0\pi R_0^2
\end{equation}

\textbf{Outlet contribution:} \\
Using polar coordinates, the area element is
\begin{equation}
    dA = 2\pi r\,dr
\end{equation}
Thus,
\begin{align*}
    \int_{A_1}(\vec{V}\cdot\hat{n})\,dA
    &=
    \int_0^{R_1}
    V_m\left(1-\frac{r^2}{R_1^2}\right)(2\pi r)\,dr \\
    &=
    2\pi V_m
    \left[
        \frac{r^2}{2}-\frac{r^4}{4R_1^2}
    \right]_0^{R_1} \\
    &=
    2\pi V_m
    \left(
        \frac{R_1^2}{2}-\frac{R_1^2}{4}
    \right) \\
    &=
    \frac{\pi}{2}V_m R_1^2
\end{align*}

\textbf{Combine both contributions:}
\begin{equation}
    0 = -V_0\pi R_0^2 + \frac{\pi}{2}V_m R_1^2
\end{equation}

Solving for the maximum outlet velocity gives
\begin{resultbox}
\begin{equation*}
    V_m = 2V_0\left(\frac{R_0}{R_1}\right)^2
\end{equation*}
\end{resultbox}